% Generated draft: requires mathematical and source review.
\documentclass[11pt]{article}
\usepackage[margin=25mm]{geometry}
\usepackage{fontspec}
\setmainfont{Latin Modern Roman}
\usepackage{amsmath,amssymb,mathtools,unicode-math}
\setmathfont{Latin Modern Math}
\usepackage{enumitem,needspace,xurl,hyperref}
\setlength{\parindent}{0pt}
\setlength{\parskip}{5pt}
\setlength{\emergencystretch}{3em}
\newcommand{\N}{\mathbb{N}}
\newcommand{\Z}{\mathbb{Z}}
\newcommand{\Q}{\mathbb{Q}}
\newcommand{\R}{\mathbb{R}}
\newcommand{\C}{\mathbb{C}}
\newcommand{\F}{\mathbb{F}}
\newcommand{\A}{\mathbb{A}}
\newcommand{\PP}{\mathbb{P}}
\newcommand{\E}{\mathbb{E}}
\newcommand{\eps}{\varepsilon}
\newcommand{\dd}{\,\mathrm{d}}
\begin{document}
\section*{UnsolvedMath ID: 30000904}
\textbf{Source problem code:} catalogue entry
\subsection{Definitions and mathematical statement}
For an infinite cardinal \(\kappa\), write \([\kappa]^\kappa=\{A\subseteq\kappa:|A|=\kappa\}\). A family \(\mathcal A\subseteq[\kappa]^\kappa\) is almost disjoint if \(|A\cap B|<\kappa\) for distinct \(A,B\in\mathcal A\). It is maximal almost disjoint if no additional member of \([\kappa]^\kappa\) can be adjoined while preserving this property; equivalently, every \(B\in[\kappa]^\kappa\) has \(|A\cap B|=\kappa\) for some \(A\in\mathcal A\).

Interpreting the source's cofinality lower bound as a restriction on the families over which the minimum is taken, define
\[
\mathfrak a_\kappa=\min\{|\mathcal A|:
\mathcal A\subseteq[\kappa]^\kappa\text{ is maximal almost disjoint},\quad
|\mathcal A|\geq\operatorname{cf}(\kappa)\},
\qquad \mathfrak a=\mathfrak a_\omega.
\]
This restriction excludes the trivial small maximal families; the intended convention at singular cardinals should be confirmed from the report.

The two separate consistency questions are whether the respective theories
\[
\mathrm{ZFC}+\bigl(\mathfrak a_{\aleph_\omega}<\mathfrak a
\ \text{and}\ \mathfrak a>\aleph_\omega\bigr)
\]
and
\[
\mathrm{ZFC}+\bigl(\mathfrak a_{\aleph_{\omega_1}}<\mathfrak a_{\aleph_1}\bigr)
\]
are consistent. Here \(\omega_1\) is the first uncountable ordinal, so the cardinal index in the second question is \(\aleph_{\omega_1}\). The extra lower bound on \(\mathfrak a\) is explicitly part of the first question only. No simultaneous consistency claim, specified forcing construction, or large-cardinal hypothesis is supplied.
\subsection{Short English statement}
Can the almost-disjointness number at \(\aleph_\omega\) be below the classical \(\mathfrak a\) while \(\mathfrak a>\aleph_\omega\)? Separately, can \(\mathfrak a_{\aleph_{\omega_1}}\) be below \(\mathfrak a_{\aleph_1}\)?
\subsection{Sources}
\begin{itemize}
\item Local UnsolvedMath registry, permanent ID \(30000904\), source code catalogue entry, ``Almost-Disjointness Numbers at Singular Cardinals.'' No individual proposer is supplied.
\item \emph{Set Theory} (2008), Oberwolfach Reports, DOI 10.4171/owr/2008/02, pp. 79--118: \url{https://doi.org/10.4171/owr/2008/02}.
\item Supplied reference, ``singular madness results'': \url{https://shelah.logic.at/files/182061/793.pdf}.
\item Supplied literature triage, checked 2026-08-21: status ``open.'' It says the report poses both inequalities as open and that a known consistency result does not decide the requested conditions, mentioning \(\mathfrak a>\aleph_\omega\). That result and its hypotheses are not reproduced. This dated assessment derives from supplied OWR curation; time-sensitive status and cited claims require independent verification.
\end{itemize}
\end{document}
